% Figure: block diagram of the sensor signal-conditioning chain, from a
% Wheatstone bridge strain sensor through an instrumentation amplifier and a
% second-order Sallen-Key anti-alias filter into an analog-to-digital
% converter. Plain TikZ boxes and arrows; no circuitikz.
\begin{figure}[htbp]
\centering
\begin{tikzpicture}[
  line width=0.8pt,
  blk/.style={draw, engblue, fill=white, minimum width=21mm, minimum height=12mm,
    align=center, font=\scriptsize},
  sig/.style={-{Latex[length=2mm]}, engblack},
  font=\small,
]
  \node[blk] (src) at (0,0) {Bridge\\sensor\\$2\,$mV/V};
  \node[blk] (ina) at (3.0,0) {Instr.\\amp\\$G=100$};
  \node[blk] (filt) at (6.0,0) {Sallen--Key\\$2$-pole\\$f_c=1\,$kHz};
  \node[blk] (adc) at (9.0,0) {ADC\\$16$-bit\\$f_s=8\,$kHz};
  \draw[sig] (src) -- (ina);
  \draw[sig] (ina) -- (filt);
  \draw[sig] (filt) -- (adc);
  % excitation arrow into the bridge
  \draw[sig, engorange] (0,1.4) -- (src.north);
  \node[font=\tiny, anchor=south, engorange] at (0,1.45) {$V_{\text{exc}}=5\,$V};
  % annotated signal levels under each interface
  \node[font=\tiny, anchor=north, enggray] at (1.5,-0.95) {$\le 10\,$mV};
  \node[font=\tiny, anchor=north, enggray] at (4.5,-0.95) {$\le 1\,$V};
  \node[font=\tiny, anchor=north, enggray] at (7.5,-0.95) {band-limited};
\end{tikzpicture}
\caption[Sensor signal-conditioning chain]{The four stages of the
signal-conditioning chain carried in the case study. A Wheatstone bridge
turns strain into a small differential voltage; an instrumentation amplifier
raises it by a factor of one hundred while rejecting the common-mode
excitation; a second-order Sallen--Key low-pass filter removes content above
the converter's Nyquist frequency; and the analog-to-digital converter
samples the band-limited result. Each stage is specified by a gain, a
bandwidth, and a noise contribution, and the chain is designed so that no
stage dominates the error budget set by the converter's least significant
bit.}
\label{fig:vol04ch09:signal-chain}
\end{figure}
